\section{Does coordination retain individual access?}\label{sec:rental-access}
The missing jobs records prevent a direct audit of its individual transitions.
We can test the same repair using the public rental GPT menus. This follow-up
reconstructs all 1,280 supported markets across the original development,
held-out and replication windows, including both supply designs. It uses
\texttt{gpt-6-luna}, not the Haiku menus in Section~\ref{sec:twonyc}.
The outcomes of these windows were already known when we designed the audit.
Their names describe the original experiment; none is an untouched test of
this follow-up.

Write $A(G)$ for a maximum-cardinality, maximum-reuse anchor applied to GPT
menus $G$. We keep every personal menu length fixed and rerun DA with the
same modeled preferences and listing priorities. Placement includes everyone.
Here coverage is the equal-market mean fraction of full-list recipients whose
partner appears on the menu. This differs from the unconditional, geographically
weighted Haiku rates above. Table~\ref{tab:rental-access} reports paired
simulation intervals adjusted jointly for eight planned contrasts within
each window and design, as specified in Appendix~\ref{app:rental-access}.

\begin{table}[ht]\centering\small
\caption{Rental GPT follow-up: $A(G)-G$, in percentage points, with 95\%
simultaneous seed-bootstrap intervals. Exits are actual replay transition
counts, not losses observed in the housing market. Appearances can repeat
households across draws.}\label{tab:rental-access}
% Generated by rental-access/rebuild.py; do not edit.
\begin{tabular}{@{}llrrr@{}}\toprule
Design & Window & $\Delta$ placement & $\Delta$ coverage & Exits / appearances\\\midrule
Ordinary & Development 2026 & $+0.69$ [-0.72, 2.10] & $-0.06$ [-1.47, 1.35] & 25/15,360 \\
Ordinary & Held-out 2026 & $+0.79$ [-1.26, 2.84] & $-0.01$ [-2.07, 2.04] & 7/10,240 \\
Ordinary & Replication 2025 & $+0.96$ [-0.15, 2.08] & $-0.13$ [-1.25, 0.98] & 25/15,360 \\
Enriched & Development 2026 & $+3.51$ [2.47, 4.55] & $-0.16$ [-1.20, 0.88] & 87/15,360 \\
Enriched & Held-out 2026 & $+4.02$ [2.83, 5.22] & $-0.33$ [-1.52, 0.86] & 37/10,240 \\
Enriched & Replication 2025 & $+4.92$ [4.03, 5.81] & $-0.25$ [-1.14, 0.64] & 78/15,360 \\
\bottomrule\end{tabular}

\end{table}

The anchor gains 3.51--4.92 placement points in enriched markets, with
positive simultaneous intervals. The gains are only 0.69--0.96 points in
ordinary markets, whose simultaneous intervals include zero. Estimated
coverage changes are negative in all six panels. None excludes zero after
the eight-contrast adjustment; this does not establish coverage equivalence.
The net gains also conceal incumbent losses: enriched panels have 87, 37 and
78 exits, and 472, 275 and 457 appearances move to a worse modeled rank,
counting loss of a seat as worse than any feasible apartment.

The coordinated benchmark again matters. Relative to anchored random,
anchored GPT's placement differences are $-0.59$, $-0.61$ and $-1.28$ points
in ordinary markets, and $+0.92$, $+0.07$ and $+1.21$ in enriched markets.
Only the last window in each design excludes zero under the simultaneous
adjustment. It retains higher reference coverage than anchored random in
all six panels. These comparisons do not support one universally best policy.

A stronger congestion control changes the conclusion about placement.
$B(G)$ rebuilds all menus to minimize squared listing exposure, then preserves
as much cached GPT rank score as that minimum permits. Against $A(G)$ it
gains 6.61--7.62 placement points in ordinary markets and 12.65--12.83 in
enriched markets, while losing 5.13--6.06 and 3.05--4.17 coverage points,
respectively. Both directions exclude zero in all six simultaneous comparisons.
This is a tradeoff between these policies, not an impossibility result.
Balancing uses the same information and menu lengths but can edit several
positions, so this is not a comparison between two one-slot repairs.
Indeed, $B(G)$ has higher placement and coverage than uncoordinated random
menus in all six panels.
A control using the simulator's observable utility scores improves both
coverage and placement over $A(G)$ in every panel with positive pointwise
intervals; it uses information beyond the cached GPT rankings.

Individual protection remains a separate requirement. Across the enriched
panels, 11 exits occur among 8,337 appearances with only one to five feasible
options. Their menus already contain every option and are unchanged by the
anchor. Other renters' edits can therefore remove access even when a person's
own menu is complete. This is the same allocation externality studied in
Chile, under a different preference model and without a field-effect claim.
Spillovers can also help: in the enriched held-out panel, balancing raises
placement by 14.41 points among renters with one to five feasible options,
whose menus remain fixed. Even this larger improvement is not individual
protection: balancing causes 171 incumbent exits overall relative to $A(G)$.
